\section{Tổng quan}

\subsection{Mô tả đề tài}

\subsubsection{Khảo sát sơ bộ}

Để hiểu rõ bài toán trước khi đề xuất giải pháp, nhóm đã thực hiện khảo sát
sơ bộ với hai nhóm đối tượng:

\textbf{Phía Giảng viên:} Qua trao đổi với \ldots\ giảng viên tại Khoa Khoa
học và Kỹ thuật Máy tính, ĐHBK TP.HCM, nhóm ghi nhận:
\begin{itemize}
    \item \textbf{\ldots\%} giảng viên phải đầu tư hơn 2 giờ để chuyển đổi
          một chương tài liệu thành nội dung kiểm tra.
    \item \textbf{\ldots\%} không có công cụ hỗ trợ, phải làm hoàn toàn thủ
          công (soạn Word, nhập Moodle).
    \item \textbf{\ldots\%} bày tỏ sẵn sàng dùng công cụ AI nếu kết quả đạt
          $\geq 70\%$ chất lượng so với làm tay.
\end{itemize}

\textbf{Phía Học viên:} Qua khảo sát \ldots\ sinh viên đại học:
\begin{itemize}
    \item \textbf{\ldots\%} cho rằng tài liệu học quá dài, khó tập trung khi
          ôn tập.
    \item \textbf{\ldots\%} muốn có bài kiểm tra nhanh sau mỗi section để
          tự đánh giá mức độ hiểu bài.
    \item \textbf{\ldots\%} đã từng dùng Quizlet hoặc flashcard để học.
\end{itemize}

\begin{figure}[H]
    \centering
    % \includegraphics[width=0.8\textwidth]{figures/survey_chart.png}
    \caption{Kết quả khảo sát nhu cầu của giảng viên và học viên.}
    \label{fig:survey}
\end{figure}

Kết quả khảo sát xác nhận rõ khoảng trống: \textbf{thiếu công cụ tự động hóa
việc tạo nội dung Micro-learning từ tài liệu sẵn có}.

\subsubsection{Mục tiêu đề tài}

\textbf{Mục tiêu tổng quát:} Xây dựng hệ thống Micro-learning thông minh có
khả năng tự động chuyển đổi tài liệu học thuật (PDF, DOCX, PPTX) thành nội
dung học tập dạng thẻ (Card) kèm bộ câu hỏi trắc nghiệm (Quiz), với sự kiểm
soát và phê duyệt của giảng viên (\textit{Human-in-the-Loop}).

\textbf{Mục tiêu cụ thể Giai đoạn 1 (Tuần 1--15):}
\begin{enumerate}
    \item Nghiên cứu và thiết kế hệ thống prompt hiệu quả cho tác vụ tóm tắt
          (Summary) và tạo câu hỏi trắc nghiệm (Quiz) từ văn bản.
    \item Xây dựng module \textbf{Text Chunking} phân đoạn văn bản theo cấu
          trúc ngữ nghĩa của tài liệu.
    \item Xây dựng \textbf{Backend AI Pipeline} hoàn chỉnh:
          Upload $\to$ Parse $\to$ Chunk $\to$ Embed $\to$ Generate.
    \item Xây dựng \textbf{giao diện web} cơ bản cho Giảng viên và Học viên.
    \item \textbf{Demo} end-to-end: upload một chương sách giáo khoa và nhận
          lại khóa học Micro-learning hoàn chỉnh trong thời gian thực.
\end{enumerate}

\subsubsection{Định hướng sản phẩm}

Sản phẩm hướng đến \textbf{nền tảng web SaaS} (Software as a Service) phục
vụ các tổ chức giáo dục đại học và trung học phổ thông tại Việt Nam. Trong
giai đoạn đầu, tập trung vào thị trường trong trường đại học.

\textbf{Mô hình sử dụng:}
\begin{itemize}
    \item Giảng viên đăng nhập, upload tài liệu, nhận kết quả AI và publish.
    \item Học viên đăng nhập, học theo Card và làm Quiz để kiểm tra kiến thức.
    \item Quản trị viên quản lý tài khoản và theo dõi thống kê hệ thống.
\end{itemize}

\subsubsection{Tiêu chí sản phẩm}

\begin{table}[H]
    \centering
    \caption{Tiêu chí chấp nhận sản phẩm Giai đoạn 1.}
    \label{tab:product-criteria}
    \begin{tabular}{p{0.5cm}p{5cm}p{3cm}p{6cm}}
        \toprule
        \textbf{\#} & \textbf{Tiêu chí} & \textbf{Ngưỡng} & \textbf{Cách đo} \\
        \midrule
        1 & Thời gian xử lý tài liệu 20 trang & $\leq$ 3 phút
          & Đo thực tế với 5 tài liệu mẫu \\
        2 & Tỷ lệ Summary đạt chất lượng & $\geq$ 70\%
          & Giảng viên đánh giá (Likert 1--5, $\geq$3 là đạt) \\
        3 & Tỷ lệ Quiz đạt chất lượng & $\geq$ 65\%
          & Tương tự trên \\
        4 & API response time (xem Card, Quiz) & $\leq$ 500ms
          & Đo với 50 concurrent users \\
        5 & SUS Score giao diện học viên & $\geq$ 68/100
          & Khảo sát 10 sinh viên \\
        6 & Coverage unit test & $\geq$ 70\%
          & \texttt{pytest --cov} \\
        \bottomrule
    \end{tabular}
\end{table}

\subsubsection{Mô hình dự kiến}

\begin{figure}[H]
    \centering
    % \includegraphics[width=0.85\textwidth]{figures/product_model.png}
    \caption{Mô hình tổng quát của sản phẩm.}
    \label{fig:product-model}
\end{figure}

\subsection{Phân công nhiệm vụ}

\subsubsection{Phân công Giai đoạn 1}

\begin{table}[H]
    \centering
    \caption{Phân công nhiệm vụ Giai đoạn 1.}
    \label{tab:task-assign-gd1}
    \begin{tabular}{p{0.5cm}p{4cm}p{2cm}p{8cm}}
        \toprule
        \textbf{\#} & \textbf{Thành viên} & \textbf{MSSV} & \textbf{Nhiệm vụ} \\
        \midrule
        1 & Nguyễn Nam Phương & 2212698
          & Nghiên cứu Prompt Engineering; thiết kế và thử nghiệm các prompt
            template; tích hợp LLM API (OpenAI/Gemini); viết chương 1 và 3. \\
        \midrule
        2 & Nguyễn Minh Triết & 2213608
          & Xây dựng module Text Chunking và Document Parsing (PDF, DOCX, PPTX);
            thiết kế kiến trúc DDD; viết chương 4 và 5. \\
        \midrule
        3 & Phùng Xương Cận & 2210348
          & Xây dựng giao diện web (Frontend Giảng viên \& Học viên); triển
            khai Docker Compose; thực hiện kiểm thử; viết chương 2 và 6. \\
        \bottomrule
    \end{tabular}
\end{table}

\subsection{Cấu trúc báo cáo}

Báo cáo được tổ chức thành 7 chương:
\begin{itemize}
    \item \textbf{Chương 1 --- Tổng quan:} Khảo sát sơ bộ, mục tiêu, tiêu
          chí sản phẩm và phân công nhiệm vụ.
    \item \textbf{Chương 2 --- Khảo sát và phân tích:} Khảo sát thực trạng
          chi tiết, phân tích giải pháp tương tự, user stories và đặc tả
          yêu cầu (BR / FR / NFR).
    \item \textbf{Chương 3 --- Cơ sở lý thuyết:} Lý thuyết LLM, Prompt
          Engineering, Text Chunking, Vector Embedding và lý giải lựa chọn
          công nghệ.
    \item \textbf{Chương 4 --- Thiết kế hệ thống:} Kiến trúc, cơ sở dữ liệu,
          giao diện (Sitemap), use case diagram và activity diagram.
    \item \textbf{Chương 5 --- Hiện thực hệ thống:} Nền tảng sử dụng, mã
          nguồn các module chính, giao diện web.
    \item \textbf{Chương 6 --- Kiểm thử hệ thống:} Kiểm thử tính năng
          (black-box), kiểm thử phi tính năng (load test, chất lượng AI).
    \item \textbf{Chương 7 --- Kết luận:} Kết quả, điểm cần cải thiện và
          định hướng Giai đoạn 2.
\end{itemize}

\subsection{Tổng kết Chương 1}

Chương 1 đã trình bày tổng quan về đề tài: bối cảnh và động cơ từ kết quả
khảo sát sơ bộ, mục tiêu cụ thể của Giai đoạn 1, tiêu chí chấp nhận sản
phẩm và phân công nhiệm vụ giữa các thành viên. Chương tiếp theo sẽ đi sâu
vào khảo sát thực trạng và phân tích yêu cầu chi tiết của hệ thống.
